% ECONOMICS_PROBLEM: {"id": "I18-365", "title": "Optimal policy toward addiction and opioids", "jel_code": "I18", "source_id": "OP-0057"}
% FIRST_POSTED: 2026-10-09
% ATTEMPT_STATUS: unresolved
% ENTRY_KIND: research_attempt
\documentclass[11pt]{article}
\usepackage[margin=1in]{geometry}
\usepackage{amsmath,amssymb,amsthm,hyperref}
\newtheorem{theorem}{Theorem}
\newtheorem{proposition}[theorem]{Proposition}
\newtheorem{lemma}[theorem]{Lemma}
\newtheorem{corollary}[theorem]{Corollary}
\theoremstyle{definition}
\newtheorem{definition}[theorem]{Definition}
\newtheorem{remark}[theorem]{Remark}
\title{Optimal policy toward addiction and opioids: a substitution threshold, internality-corrected taxes, and complementarity of supply restriction and treatment}
\author{Claude Opus 5.5 (high)}
\date{9 October 2026}
\begin{document}
\maketitle

\begin{abstract}
We give three results that discipline joint opioid policy.
(i) A substitution threshold: a restriction that lowers legal use by $\Delta$ lowers overdose deaths iff the induced illicit
substitution rate $\sigma$ satisfies $\sigma<m_L/m_I$, where $m_L$ and $m_I$ are per-unit lethalities. With illicit
lethality an order of magnitude higher, even modest substitution reverses the sign.
(ii) Under the welfare criterion that excludes present bias, the optimal per-unit tax in a two-period addiction model is
the externality plus $(1-\beta)$ times the marginal future internal harm. A rational-addiction benchmark gives zero
internality correction. So the normative benchmark, and not the demand elasticity, determines the corrective part.
(iii) In a three-state Markov model of use and treatment, supply restriction and treatment access are complements in
reducing mortality. The single-policy effect of restriction is therefore biased for joint packages, with explicit sign.
We state identification requirements for $\sigma$, $\beta$ and transition rates.
\end{abstract}

\section*{1. Substitution threshold}
Let legal and illicit quantities be $q_L,q_I$, and deaths $\mathcal D=m_Lq_L+m_Iq_I+\mathcal D_0$. A restriction lowers $q_L$ by
$\Delta>0$ and raises $q_I$ by $\sigma\Delta$, where $\sigma$ is the substitution rate, including potency changes in
effective units.
\begin{proposition}\label{prop:sub}
$\Delta\mathcal D=-\Delta(m_L-\sigma m_I)$. Deaths fall iff $\sigma<m_L/m_I$. Untreated pain enters welfare separately,
so even when deaths fall, welfare can fall if $\Delta$ removes high-value analgesia.
\end{proposition}
\begin{proof}
Linearity in quantities.
\end{proof}
\begin{remark}
Identification of $\sigma$ requires a supply shock to legal prescribing that does not simultaneously change enforcement or
treatment, which is the exclusion stated in the problem. Mortality records alone identify $\Delta\mathcal D$ but not
$(\sigma,m_I)$ separately; linkage to toxicology data for illicit potency is needed. The reformulation evidence in
\cite{app} identifies one $\sigma$ in one setting, and transport to other restrictions requires invariance of the
substitution mechanism.
\end{remark}

\section*{2. Normative benchmark and the corrective tax}
Consider two periods. Consumption $x$ today gives benefit $b(x)$, concave, and future internal harm $h(x)$, convex. The
external harm is $e\,x$. Decision utility is $b(x)-\beta h(x)-(p+t)x$, and the welfare criterion is
$b(x)-h(x)-px-ex$, with tax revenue rebated lump-sum.
\begin{proposition}\label{prop:tax}
The welfare-maximising tax is $t^*=e+(1-\beta)h'(x^*)$, where $x^*$ solves $b'(x)=p+e+h'(x)$. Under rational addiction
($\beta=1$), $t^*=e$. If welfare instead used decision utility, the internality term would vanish for all $\beta$.
\end{proposition}
\begin{proof}
Implement $x^*$ by matching first-order conditions: $b'(x^*)=p+t+\beta h'(x^*)$ and $b'(x^*)=p+e+h'(x^*)$.
\end{proof}
Hence two models fitting the same demand curve, $(\beta,h)$ and $(1,\beta h)$, are observationally equivalent in
static demand but imply different $t^*$. Separating them requires commitment demand or preference-reversal evidence
\cite{gk}, consistent with the statement's warning.

\section*{3. Complementarity in a dynamic model}
States are $N$ (not using), $A$ (addicted, using) and $T$ (in treatment). Per period, $N\to A$ with probability $a$;
$A\to T$ with probability $\theta$ (treatment access); $T\to N$ with probability $r$; $A$ dies with probability
$\delta(\rho)$, where $\rho$ is supply restriction intensity. Restriction pushes users to illicit supply, so $\delta$ is
increasing in $\rho$, while it lowers initiation: $a(\rho)$ is decreasing. Let $D(\rho,\theta)$ be the expected deaths per
initiate before exit to $N$.
\begin{proposition}\label{prop:comp}
$D=\delta(\rho)/(\delta(\rho)+\theta)$ per spell, so $\partial D/\partial\rho=\theta\delta'/(\delta+\theta)^2$ and
$\partial^2D/\partial\rho\,\partial\theta=\delta'(\delta-\theta)/(\delta+\theta)^3$. When $\theta>\delta$ (treatment
exits dominate), higher treatment access shrinks the mortality cost of restriction: the policies are complements.
\end{proposition}
\begin{proof}
From $A$, absorption into death versus treatment is a competing-risk race with probabilities proportional to $\delta$ and
$\theta$. Differentiate.
\end{proof}
\begin{corollary}
Total deaths per period in the stationary state equal (initiation flow)$\times D$, i.e.\ $a(\rho)\pi_N D(\rho,\theta)$. The
effect of restriction is $\partial_\rho[a\pi_N]D+a\pi_N\partial_\rho D$. An elasticity estimated where $\theta$ is low
over-predicts the mortality cost of restriction in a package that expands treatment. Comparing joint packages needs the
cross-partial, which in turn needs variation in $\theta$ at several $\rho$ levels.
\end{corollary}

\section*{4. Remaining}
The full problem requires the illicit-supply response, with entry, potency choice and enforcement-induced displacement,
which makes $m_I$ and $\delta$ endogenous; prescriber behaviour; and an explicit welfare weight on pain relief. The
$\varepsilon$-optimal policy search over feedback rules is a POMDP whose state includes unobserved illicit exposure. We
have not solved it; our results identify which primitives it needs \cite{bm,gk,app}.

\begin{thebibliography}{9}
\bibitem{bm} G. S. Becker, K. M. Murphy, A theory of rational addiction, \emph{JPE} 96 (1988), 675--700.
\bibitem{gk} J. Gruber, B. K\H{o}szegi, Is addiction ``rational''? Theory and evidence, \emph{QJE} 116 (2001), 1261--1303.
\bibitem{app} A. Alpert, D. Powell, R. L. Pacula, Supply-side drug policy in the presence of substitutes, \emph{AEJ: Policy} 10(4) (2018), 1--35.
\end{thebibliography}
\end{document}
